% !TeX root = ../../Book3.tex
\section{Dirichlet 级数与 Euler 乘积}
\label{b3:sec:2601}

本章把上一章的收敛级数扩展为全平面的亚纯函数。首先保留它与素数的直接关系；随后用高斯的参数反演得到在原级数收敛域之外仍有意义的表达式。

\subsection{在 \texorpdfstring{\(\operatorname{Re}s>1\)}{Re(s)>1} 上的收敛性}
\label{b3:subsec:260101}

\begin{definition}\label{b3:def:26-zeta}
在 \(\operatorname{Re}s>1\) 上定义\term[riemann-zeta-hanshu]{Riemann zeta 函数}[Riemann zeta function]
\(\zeta(s)\coloneqq\sum_{n\ge1}n^{-s}\)。
\end{definition}
\symidx{\(\zeta(s)\)：Riemann zeta 函数}
对紧集 \(\operatorname{Re}s\ge1+\delta\)，由 \(\sum n^{-1-\delta}<\infty\) 得局部一致绝对收敛，故全纯；任意固定次导数可逐项计算，因为 \((\log n)^k\) 被小幂 \(n^{\delta/2}\) 控制。在正实轴，\(\zeta(\sigma)>1\)，且 \(\zeta(\sigma)\to1\) 当 \(\sigma\to\infty\)。

\subsection{Euler 乘积}
\label{b3:subsec:260102}

\cref{b3:thm:25-euler-product}给 \(\zeta(s)=\prod_p(1-p^{-s})^{-1}\)。更强的是前章已经建立收敛全纯对数
\(\log\zeta(s)=\sum_{p,r\ge1}p^{-rs}/r\)，从而在这个半平面无零点。若把乘积形式延伸到 \(\operatorname{Re}s\le1\)，它通常不再收敛；后面的解析延拓不能被称为该无穷乘积在原意义下的值。

\ProseParagraphBreak

绝对收敛还给出不依赖虚部的上下界。由原级数及 Möbius 逆级数，
\[
 |\zeta(\sigma+\mathrm it)|\le\zeta(\sigma),\qquad
 \left|\frac1{\zeta(\sigma+\mathrm it)}\right|
 \le\sum_{n\ge1}\frac{|\mu(n)|}{n^\sigma}\le\zeta(\sigma),
 \quad\sigma>1.
\]
所以在每条固定的直线 \(\sigma>1\) 上，函数既有上界，也与零保持正距离。不过当 \(\sigma\downarrow1\) 时，这个下界趋于零，因而不能排除边界零点。无零点区域是否能够扩大，必须从新的信息判断。

\subsection{素数与 \texorpdfstring{\(\zeta\)}{\AnalysisPdfZeta} 函数}
\label{b3:subsec:260103}

\(\zeta\) 将所有整数编码，而其对数只在素数幂上有系数。沿实轴 \(\sigma\downarrow1\)，\(\zeta(\sigma)\) 的发散迫使素数无限多，也迫使 \(\sum_p1/p\) 发散。这类结论仍比素数定理弱：它们不能确定 \(\pi(x)\) 的渐近常数，后者需要控制函数在边界直线上的零点。
